\section{強化学習}

強化学習（reinforcement learning）は，環境との相互作用を通じて意思決定規則を学習し，長期的な累積報酬を最大化するための理論である．教師あり学習では入力と正解の対応が与えられ，教師なし学習では観測データの潜在構造の抽出が中心となるのに対し，強化学習では各時刻の行動に対して正解が即座に与えられるとは限らない．ある行動の評価は，その後に続く状態遷移と報酬列を通じて初めて定まり，その意味で遅延報酬を本質に含んでいる．

本章の目標は，強化学習を単なるアルゴリズムの集合としてではなく，確率的環境のもとで価値関数または方策を学習し，累積報酬を最適化する統一理論として理解することである．標準的な教科書的整理としては \citet{SuttonBarto2018} が広く参照されるが，本章ではそれを踏まえつつ，交通・都市工学の逐次意思決定問題にも接続しやすい形で議論を組み立てる．まず歴史的流れから問題意識を整理し，次にマルコフ決定過程によって標準的定式化を与える．そのうえで，多腕バンディット問題によって探索と活用の本質を切り出し，動的計画法，価値ベース手法，方策勾配法へと進む．最後に，標準的な設定だけでは捉えきれない発展的話題を概観する．

\subsection{強化学習の歴史}

強化学習の歴史を学ぶ必要があるのは，この分野が単に高性能な手法を寄せ集めて成立したのではなく，逐次意思決定に固有の困難に対する応答として発展してきたからである．信号制御，車両経路選択，需要応答交通の運行計画のような問題では，現在の行動が将来の混雑や待ち時間を変化させ，評価は複数時刻にまたがって現れる．したがって，入力と正解の対応をそのまま学べばよい教師あり学習とは異なり，行動の価値を試行錯誤から推定しなければならない．

このとき難しさの中心にあるのが，遅延報酬，探索と活用のトレードオフ，そして方策依存のデータ生成である．第一に，最終的な成果にどの行動が寄与したかを識別する信用割当ての問題が生じる．第二に，より良い行動を見つけるためには未知の選択肢を試す探索が必要だが，探索を増やせば当面の報酬は失われる．第三に，データは受動的に与えられるのではなく，学習者自身の行動によって生成される．これらの特徴が，強化学習に独自の理論を要請する．

歴史的源流は，Bellman による動的計画法と最適制御にある\cite{Bellman1957,Puterman1994}．ここで与えられた基本原理は，長期最適化問題を現在の価値と将来の価値の再帰関係として表すベルマン方程式であった．ただし，この枠組みは環境モデル，すなわち状態遷移確率と報酬関数が既知であることを前提としていた．理論的には明快であっても，未知環境から経験を通じて学ぶという強化学習本来の状況には，そのままでは適用できなかった．

この不足を埋めたのが時間差分学習（temporal-difference learning）である．\citet{Sutton1988}は，将来の完全な帰結を待たず，次時刻の価値推定を足場にして現在の価値を更新するブートストラップの考え方を定式化した．これにより，モンテカルロ法のようにエピソード終了まで待たなくても逐次的に価値を学習できるようになった．ベルマン方程式は，解析のための式であるだけでなく，学習更新式の源泉になったのである．

さらに，\citet{WatkinsDayan1992}による Q 学習は，行動価値関数を直接学ぶことで，環境モデルなしに最適方策へ近づく道を与えた．ここで重要なのは，価値推定と制御の問題が統合されたことである．TD 学習が与えられた方策の価値推定を中心にしていたのに対し，Q 学習は最大演算を通じて最適方策そのものを目指した．

状態空間が高次元化すると，表形式の価値関数では表現能力が不足する．この制約を緩和したのが深層強化学習であり，とくに DQN は画像入力から行動価値を学ぶ代表的成功例となった\cite{Mnih2015}．AlphaGo 以降は，価値学習，方策学習，探索，計画を組み合わせる深層強化学習が急速に発展した\cite{Silver2016}．ただし，この発展は流行紹介として理解するべきではない．本質は，既知モデルから未知モデルへ，表形式から関数近似へ，価値推定から方策最適化へと理論の焦点が拡張された点にある．

同時に，価値ベース手法だけでは連続行動空間や確率的方策の直接最適化に限界があることも明らかになった．\citet{Williams1992}に始まる方策勾配法は，方策を直接最適化対象とする考え方を与え，後に Actor-Critic や PPO のような実用的手法へ発展した\cite{Schulman2015,Schulman2017}．この歴史を貫いているのは，強化学習を「環境との相互作用のもとで累積報酬を最大化する確率的意思決定問題」としてどう定式化し，どう近似的に解くかという問いである．次節では，その標準的枠組みとしてマルコフ決定過程を導入する．

\subsection{マルコフ決定過程}

強化学習を厳密に論じるには，何が状態であり，何が行動であり，何が確率的に変化し，何を最大化したいのかを明示する必要がある．そのための標準的枠組みがマルコフ決定過程（Markov decision process; MDP）である．MDP によって，遅延報酬，価値関数，方策最適化という本章の主要概念を一つの数理的言語で記述できる．

\paragraph*{定義}
MDP は，状態空間 $\mathcal{S}$，行動空間 $\mathcal{A}$，遷移確率 $P(s' \mid s,a)$，報酬関数 $r(s,a,s')$，割引率 $\gamma$ からなる組
\begin{equation}
(\mathcal{S}, \mathcal{A}, P, r, \gamma)
\end{equation}
として与えられる．ここで $s \in \mathcal{S}$ は現在状態，$a \in \mathcal{A}$ は現在行動，$s' \in \mathcal{S}$ は次状態を表す．遷移確率 $P(s' \mid s,a)$ は，状態 $s$ で行動 $a$ を選んだときに次状態が $s'$ になる条件付き確率であり，報酬関数 $r(s,a,s')$ はその遷移に対応する即時報酬の期待値である．割引率 $\gamma \in [0,1)$ は，将来報酬をどの程度現在価値として重視するかを表す係数である．

時刻 $t$ に状態 $s_t$ を観測し，行動 $a_t$ を選び，報酬 $r_t$ を受け取って次状態 $s_{t+1}$ へ移るとする．このとき，時刻 $t$ 以降の累積割引報酬，すなわちリターン $G_t$ を
\begin{equation}
G_t = \sum_{k=0}^{\infty} \gamma^k r_{t+k}
\end{equation}
と定義する．添字 $k$ は現在から何ステップ先の報酬かを表し，$\gamma^k$ はその報酬に掛かる割引である．有限ホライズンではこの和は終端時刻までの有限和となり，無限ホライズンでは $\gamma < 1$ によって収束が保証される．

MDP の核心はマルコフ性にある．これは，次状態と次報酬の条件付き分布が，過去全体ではなく現在状態と現在行動だけで決まるという仮定であり，
\begin{equation}
\mathbb{P}(s_{t+1},r_t \mid s_0,a_0,\ldots,s_t,a_t)
=
\mathbb{P}(s_{t+1},r_t \mid s_t,a_t)
\end{equation}
と書ける．この仮定が重要なのは，現在状態が将来予測に必要な情報を十分に要約しているとみなせるからである．その結果，長期最適化問題を再帰式として記述できる．

\paragraph*{定義}
方策（policy）$\pi(a \mid s)$ は，状態 $s$ のもとで行動 $a$ を選ぶ条件付き確率である．決定論的方策では各状態に一つの行動が対応し，確率的方策では複数の行動に確率が割り当てられる．方策 $\pi$ が与えられたとき，状態価値関数（state-value function）$V^\pi(s)$ と行動価値関数（action-value function）$Q^\pi(s,a)$ を
\begin{align}
V^\pi(s) &= \mathbb{E}_\pi[G_t \mid s_t=s], \\
Q^\pi(s,a) &= \mathbb{E}_\pi[G_t \mid s_t=s, a_t=a]
\end{align}
と定義する．$\mathbb{E}_\pi[\cdot]$ は，方策 $\pi$ と遷移確率 $P$ のもとでの期待値である．$V^\pi(s)$ は状態 $s$ にいることの長期的な良さを，$Q^\pi(s,a)$ は状態 $s$ で行動 $a$ を取ることの長期的な良さを表す．

これらの価値関数は，ベルマン期待方程式を満たす．状態価値関数については
\begin{align}
V^\pi(s)
&=
\sum_{a \in \mathcal{A}} \pi(a \mid s)
\sum_{s' \in \mathcal{S}} P(s' \mid s,a)
\left[
r(s,a,s') + \gamma V^\pi(s')
\right]
\end{align}
が成り立つ．同様に，行動価値関数は
\begin{align}
Q^\pi(s,a)
&=
\sum_{s' \in \mathcal{S}} P(s' \mid s,a)
\left[
r(s,a,s') + \gamma \sum_{a' \in \mathcal{A}} \pi(a' \mid s') Q^\pi(s',a')
\right]
\end{align}
を満たす．ここで $a' \in \mathcal{A}$ は次状態 $s'$ における次の行動である．価値が即時報酬と将来価値の和として再帰的に分解されることが，強化学習の理論的出発点である．

最適方策 $\pi^\ast$ は，全状態で価値を最大化する方策として定義される．対応する最適状態価値関数 $V^\ast(s)$ と最適行動価値関数 $Q^\ast(s,a)$ は
\begin{align}
V^\ast(s) &= \max_\pi V^\pi(s), \\
Q^\ast(s,a) &= \max_\pi Q^\pi(s,a)
\end{align}
であり，ベルマン最適方程式
\begin{align}
V^\ast(s)
&=
\max_{a \in \mathcal{A}}
\sum_{s' \in \mathcal{S}} P(s' \mid s,a)
\left[
r(s,a,s') + \gamma V^\ast(s')
\right], \\
Q^\ast(s,a)
&=
\sum_{s' \in \mathcal{S}} P(s' \mid s,a)
\left[
r(s,a,s') + \gamma \max_{a' \in \mathcal{A}} Q^\ast(s',a')
\right]
\end{align}
を満たす．最大演算は，将来にわたって最も価値の高い行動を選ぶという最適制御の原理を表している．

有限ホライズンと無限ホライズン，エピソード型と継続型の区別も重要である．有限ホライズンでは終端時刻が存在し，無限ホライズンでは割引率が重要になる．エピソード型では終端状態に達すると課題が終了するのに対し，継続型ではタスクが終わらないため，交通流制御のような問題では長期的安定性をどう定義するかが論点になる．次節では，状態遷移という複雑さを一度外し，探索と活用の本質だけを切り出した多腕バンディット問題を考える．

\subsection{多腕バンディット問題}\label{subsec:bandit}

MDP は強化学習の標準形であるが，状態遷移まで含めると探索と活用の本質が見えにくくなることがある．そこで，多腕バンディット問題（multi-armed bandit problem）を導入する．これは状態遷移を持たず，未知の行動価値を試行錯誤によって学ぶという構造だけを残した問題であり，強化学習の基礎概念を理解するうえで有用である\cite{Auer2002,LattimoreSzepesvari2020}．

\paragraph*{定義}
$K$ 本の腕を持つバンディットを考える．各腕 $i \in \{1,\ldots,K\}$ を選ぶと，未知分布に従う報酬が得られ，その期待値を $\mu_i$ とする．時刻 $t$ に選ばれる腕を $A_t$，得られる報酬を $r_t$ と書く．この問題には状態遷移がなく，意思決定は「次にどの腕を引くか」だけである．最適腕の期待報酬を $\mu^\ast=\max_i \mu_i$ とすると，目標は時刻 $T$ までの累積期待報酬を最大化することである．

このとき，性能評価には後悔（regret）がよく用いられる．累積後悔 $R_T$ を
\begin{equation}
R_T
=
T \mu^\ast - \mathbb{E}\left[\sum_{t=1}^{T} r_t\right]
\end{equation}
と定義すると，これは「常に最適腕だけを引いた場合に比べてどれだけ損をしたか」を表す．強化学習全体では累積報酬最大化が中心目的であるが，バンディットでは後悔最小化の形で理論解析が明確になる．

最も単純な戦略の一つが $\epsilon$-greedy である．これは確率 $1-\epsilon$ で現在までに推定された最良の腕を選び，確率 $\epsilon$ で一様ランダムに探索する方法である．ここで $\epsilon \in [0,1]$ は探索率を表す．直感的には，既知の有望な腕を主に活用しつつ，未知の腕をときどき試すことで誤った早期収束を防ぐ．ただし，不確実性の大きさを腕ごとに使い分けないため，探索の効率は高くない．

これに対して，上側信頼限界法（upper confidence bound; UCB）は，推定平均と不確かさの両方を用いる\cite{Auer2002}．腕 $i$ の時刻 $t-1$ までの経験平均を $\hat{\mu}_i(t-1)$，その腕を選んだ回数を $N_i(t-1)$ とすると，典型的には
\begin{equation}
A_t
=
\mathop{\mathrm{arg\,max}}_i
\left[
\hat{\mu}_i(t-1)
+
c \sqrt{\frac{\log t}{N_i(t-1)}}
\right]
\end{equation}
で腕を選ぶ．ここで $c>0$ は探索の強さを調整する定数である．右辺第二項は未観測性に対応し，試行回数が少ない腕ほど大きくなる．したがって UCB は，「現在よさそうな腕」だけでなく「まだ十分試していないので本当は良いかもしれない腕」も積極的に選ぶ．

もう一つの代表的方法が Thompson sampling である\cite{Thompson1933}．これは，各腕の平均報酬に関する事後分布を持ち，そこから仮想的な平均報酬 $\tilde{\mu}_i$ をサンプルして
\begin{equation}
A_t = \mathop{\mathrm{arg\,max}}_i \tilde{\mu}_i
\end{equation}
とする方法である．たとえばベルヌーイ報酬では，成功確率にベータ分布を置くことで計算が容易になる．Thompson sampling の特徴は，探索が外的なランダム化としてではなく，パラメータ不確実性のベイズ的表現から自然に導かれる点にある．

バンディット問題が MDP より簡単なのは，状態遷移がないため，長期的な信用割当てを考える必要がないからである．一方で，本質的に共通している点も明確である．すなわち，未知の環境で累積報酬を最大化するには，不確実性のもとで探索と活用を両立させなければならない．この意味で，多腕バンディットは強化学習の縮約版である．次節では，再び状態遷移を導入し，ただし環境モデルが既知である場合の計画法として動的計画法を扱う．

\subsection{動的計画法}

動的計画法（dynamic programming）が必要になるのは，環境モデルが既知であれば，経験から学ぶ前に理論上の最適解を計算できるからである．これは現実にはしばしば非現実的な仮定であるが，強化学習アルゴリズムの基準点としてきわめて重要である．モデル未知の場合の TD 学習や Q 学習は，多くの場合この枠組みをサンプル近似したものと理解できる．

\paragraph*{定義}
方策 $\pi$ に対するベルマン作用素 $T^\pi$ を
\begin{equation}
(T^\pi V)(s)
=
\sum_{a \in \mathcal{A}} \pi(a \mid s)
\sum_{s' \in \mathcal{S}} P(s' \mid s,a)
\left[
r(s,a,s') + \gamma V(s')
\right]
\end{equation}
と定義する．ここで $V$ は任意の状態価値関数である．同様に，最適ベルマン作用素 $T_\ast$ を
\begin{equation}
(T_\ast V)(s)
=
\max_{a \in \mathcal{A}}
\sum_{s' \in \mathcal{S}} P(s' \mid s,a)
\left[
r(s,a,s') + \gamma V(s')
\right]
\end{equation}
と定義する．ベルマン期待方程式は $V^\pi = T^\pi V^\pi$，ベルマン最適方程式は $V^\ast = T_\ast V^\ast$ という不動点方程式になる．

\paragraph*{アルゴリズム}
方策評価（policy evaluation）は，固定された方策 $\pi$ の価値関数 $V^\pi$ を求める手続きである．初期関数 $V_0$ を任意に取り，
\begin{equation}
V_{k+1} = T^\pi V_k
\end{equation}
と反復すると，$V_k$ は $V^\pi$ に近づく．直感的には，価値は終端報酬や高報酬状態から逆向きに伝播する．まず 1 ステップ先の価値が更新され，その情報がさらに前の状態へ波及していくのである．

方策改善（policy improvement）は，現在の価値関数にもとづいてより良い行動を選ぶ段階である．すなわち，各状態 $s$ について
\begin{equation}
\pi'(s)
\in
\mathop{\mathrm{arg\,max}}_{a \in \mathcal{A}}
\sum_{s' \in \mathcal{S}} P(s' \mid s,a)
\left[
r(s,a,s') + \gamma V^\pi(s')
\right]
\end{equation}
とすれば，新しい方策 $\pi'$ は少なくとも $\pi$ と同程度，通常はそれ以上の価値を持つ．これは，現在の価値評価に照らして各状態で最善の 1 ステップ行動を選べば，方策全体が改善することを意味している．

方策反復（policy iteration）は，方策評価と方策改善を交互に行う方法である．
\begin{enumerate}
\item 現在の方策 $\pi_k$ を評価して $V^{\pi_k}$ を求める．
\item $V^{\pi_k}$ に対して貪欲化を行い，新しい方策 $\pi_{k+1}$ を得る．
\item 方策が変化しなくなるまで反復する．
\end{enumerate}

価値反復（value iteration）は，方策評価を完全に解かずに最適ベルマン作用素を直接適用する方法であり，
\begin{equation}
V_{k+1} = T_\ast V_k
\end{equation}
を反復する．その後，各状態で価値を最大化する行動を選べば最適方策が得られる．価値反復は「いま分かっている将来価値に照らして最善の 1 ステップ選択」を毎回価値へ織り込む方法であり，計画と改善が一体化されている．

\paragraph*{定理}
ベルマン作用素が収束をもたらす理由は，割引率 $\gamma < 1$ のもとで縮小写像になるからである．最大ノルム $\|V-W\|_\infty = \max_{s \in \mathcal{S}} |V(s)-W(s)|$ を用いると，
\begin{align}
\|T^\pi V - T^\pi W\|_\infty &\le \gamma \|V-W\|_\infty, \\
\|T_\ast V - T_\ast W\|_\infty &\le \gamma \|V-W\|_\infty
\end{align}
が成り立つ．誤差が反復ごとに少なくとも $\gamma$ 倍に縮むため，不動点に向かって収束する．直感的には，遠い将来の影響が割引によって弱まり，誤差が増幅しないことを意味する．

動的計画法は理論上の理想解法である一方，遷移確率と報酬関数が既知でなければ使えない．そのため，モデル未知の状況では期待値計算を実際の観測遷移で置き換える必要がある．次節では，その橋渡しとして TD 学習から Q 学習へ至る価値ベース手法を扱う．

\subsection{価値ベース手法：Q学習, Deep-Q学習}

現実の強化学習では，環境モデルが未知であることが多い．したがって，動的計画法のように遷移確率の期待値を厳密に計算することはできず，経験から価値を学習しなければならない．価値ベース手法は，価値関数を学ぶことを通じて最適方策を間接的に得る方法であり，モデル未知環境における強化学習の中心的流儀である．

出発点となるのは TD 学習である．状態価値関数 $V$ に対する TD(0) の更新は
\begin{equation}
V(s_t)
\leftarrow
V(s_t) + \alpha
\left[
r_t + \gamma V(s_{t+1}) - V(s_t)
\right]
\end{equation}
で与えられる．ここで $\alpha \in (0,1]$ は学習率である．角括弧内の量
\begin{equation}
\delta_t = r_t + \gamma V(s_{t+1}) - V(s_t)
\end{equation}
を TD 誤差という．$\delta_t$ は，現在の価値予測が 1 ステップ先の観測を踏まえてどれだけ過小または過大であったかを示す量であり，価値修正の方向と大きさを与える信号である．

制御問題では，状態そのものの価値だけでなく，各行動の価値を知る必要がある．そこで行動価値関数 $Q(s,a)$ を学習する．Q 学習は
\begin{equation}
Q(s_t,a_t) \leftarrow Q(s_t,a_t) + \alpha \left[r_t + \gamma \max_{a'} Q(s_{t+1},a') - Q(s_t,a_t)\right]
\end{equation}
という更新式で定義される\cite{WatkinsDayan1992}．ここで $a' \in \mathcal{A}$ は次状態 $s_{t+1}$ における候補行動であり，$\max_{a'} Q(s_{t+1},a')$ は将来に最善行動を選ぶと仮定した価値である．この最大演算を通じて，Q 学習は最適方策へ向かう．

Q 学習の重要な特徴は off-policy 学習である点にある．すなわち，実際にデータを収集する行動方策と，更新式の内部で想定している目標方策が一致していなくてもよい．たとえば $\epsilon$-greedy によって探索を行いながら，更新時には貪欲方策を想定できる．これに対し SARSA は
\begin{equation}
Q(s_t,a_t)
\leftarrow
Q(s_t,a_t) + \alpha
\left[
r_t + \gamma Q(s_{t+1},a_{t+1}) - Q(s_t,a_t)
\right]
\end{equation}
と，実際に選ばれた次行動 $a_{t+1}$ を用いて更新する on-policy 手法である．Q 学習が将来の最良行動に基づく強気の更新であるのに対し，SARSA は探索を含んだ現在の方策をそのまま評価する．

表形式の Q 学習は，状態空間と行動空間が小さいときには有効であるが，高次元状態では破綻する．画像入力や多数の交通状態特徴量のような状況では，各 $(s,a)$ に対して表を持つことができない．このとき必要になるのが関数近似であり，Deep Q-Network（DQN）はニューラルネットワーク $Q(s,a;\theta)$ によって行動価値関数を近似する\cite{Mnih2015}．

DQN では，TD ターゲットを
\begin{equation}
y_t = r_t + \gamma \max_{a'} Q(s_{t+1},a';\theta^-)
\end{equation}
と置き，損失関数
\begin{equation}
L_t(\theta)
=
\left[
y_t - Q(s_t,a_t;\theta)
\right]^2
\end{equation}
を最小化する．ここで $\theta$ は学習中のネットワークの重み，$\theta^-$ はターゲットネットワークの重みである．もし右辺のターゲットまで同じ重みで毎回更新すると，目標そのものが激しく動くため学習が不安定化する．ターゲットネットワークはこの問題を緩和し，更新先を一時的に固定する役割を果たす．

もう一つの重要要素が経験再生（experience replay）である．これは，過去の遷移 $(s_t,a_t,r_t,s_{t+1})$ をバッファに蓄積し，そこから無作為にミニバッチを抽出して学習する仕組みである．経験再生が必要なのは，連続する遷移が強く相関しているため，そのまま学習すると勾配推定が不安定になるからである．無作為抽出により相関を弱められ，同じ経験を複数回再利用できるのでデータ効率も向上する．DQN の安定化は，この経験再生とターゲットネットワークに大きく依存している．

その後，Double DQN は最大演算に由来する過大評価を軽減し，dueling network は状態価値とアドバンテージを分離して表現することで学習効率を高めた．しかし中心的な論点は，TD 学習の再帰構造を深層関数近似とどう両立させるかにある．価値ベース手法の長所は，離散行動空間では各行動の価値を直接比較できること，ベルマン方程式との結びつきが明瞭であることである．一方で，連続行動空間では $\max_{a'} Q(s_{t+1},a')$ の計算が難しく，確率的方策の直接制御や制約付き最適化にも不向きである．この限界が，次節で扱う方策勾配法の必要性を生む．

\subsection{方策勾配法：Actor-Critic, PPO}

価値ベース手法の限界が明確になるのは，連続行動空間や確率的方策の直接最適化を考えるときである．ロボット制御のように行動が実数ベクトルで表される場合，離散行動に対する最大化は自然ではない．また，探索を方策の分布そのものに埋め込みたい場合には，価値関数を介した間接的な最適化より，方策を直接最適化する方が適切である．この問題意識から方策勾配法が導入される．

パラメータ $\theta$ を持つ方策を $\pi_\theta(a \mid s)$ とし，期待リターンを
\begin{equation}
J(\theta) = \mathbb{E}_{\pi_\theta}[G_0]
\end{equation}
と定義する．ここで $\mathbb{E}_{\pi_\theta}[\cdot]$ は，方策 $\pi_\theta$ と環境ダイナミクスに従って生成される軌道に関する期待値である．方策勾配法の目標は，$J(\theta)$ を最大化するように $\theta$ を更新することである．

\paragraph*{定理}
方策勾配の基本形は
\begin{equation}
\nabla_\theta J(\theta) = \mathbb{E}\left[\nabla_\theta \log \pi_\theta(a_t|s_t)\, G_t\right]
\end{equation}
で与えられる\cite{Williams1992}．ここで $\nabla_\theta$ はパラメータ $\theta$ に関する勾配，$\log \pi_\theta(a_t \mid s_t)$ は選択された行動の対数確率，$G_t$ は時刻 $t$ 以降のリターンである．この式は，良い帰結をもたらした行動の確率を高め，悪い帰結をもたらした行動の確率を下げるという学習規則を正当化する．

REINFORCE はこの勾配を標本平均で近似する最も基本的なアルゴリズムであり，
\begin{equation}
\theta \leftarrow \theta + \alpha \sum_t \nabla_\theta \log \pi_\theta(a_t \mid s_t) G_t
\end{equation}
と更新する．ここで $\alpha$ は学習率である．REINFORCE は不偏推定量を与えるが，$G_t$ の分散が大きく，学習が不安定になりやすい．とくにホライズンが長い問題では，分散の削減が不可欠になる．

そこでベースライン $b(s_t)$ を導入し，
\begin{equation}
\nabla_\theta J(\theta)
=
\mathbb{E}\left[
\nabla_\theta \log \pi_\theta(a_t \mid s_t)
\left(G_t - b(s_t)\right)
\right]
\end{equation}
と書き換える．$b(s_t)$ が状態のみに依存する限り，勾配の期待値は変わらず，分散だけを減らせる．とくに $b(s_t)=V^\pi(s_t)$ と置くと，アドバンテージ関数
\begin{equation}
A^\pi(s_t,a_t) = Q^\pi(s_t,a_t) - V^\pi(s_t)
\end{equation}
が現れる．これは，状態 $s_t$ において行動 $a_t$ が平均的行動よりどれだけ良かったかを表す量である．

Actor-Critic はこの考え方を実装する枠組みである．Actor は方策 $\pi_\theta(a \mid s)$ を表し，Critic は価値関数，典型的には $V(s;w)$ を推定する．ここで $w$ は Critic のパラメータである．Critic は TD 学習により
\begin{equation}
\delta_t = r_t + \gamma V(s_{t+1};w) - V(s_t;w)
\end{equation}
を計算する．この TD 誤差 $\delta_t$ は，1 ステップのアドバンテージ推定として機能する．したがって Actor は
\begin{equation}
\theta \leftarrow \theta + \alpha \, \delta_t \, \nabla_\theta \log \pi_\theta(a_t \mid s_t)
\end{equation}
と更新できる．$\delta_t$ が正なら行動確率を増やし，負なら減らすので，Critic の価値推定が Actor の改善方向を与えることになる．

実際の深層強化学習では，方策更新が大きすぎると性能が急激に悪化することがある．TRPO は新旧方策の距離に制約を課すことでこの問題に対処したが\cite{Schulman2015}，計算は比較的重い．これに対し PPO は，同じ問題意識をより簡潔な目的関数で扱う\cite{Schulman2017}．

PPO では，古い方策を $\pi_{\theta_{\mathrm{old}}}$，新しい方策を $\pi_\theta$ とし，確率比
\begin{equation}
r_t(\theta)
=
\frac{\pi_\theta(a_t \mid s_t)}{\pi_{\theta_{\mathrm{old}}}(a_t \mid s_t)}
\end{equation}
を考える．さらに $A_t$ を推定アドバンテージとすると，clip された目的関数は
\begin{equation}
L^{\mathrm{clip}}(\theta)
=
\mathbb{E}
\left[
\min\left(
r_t(\theta)A_t,\,
\mathrm{clip}\left(r_t(\theta),1-\epsilon,1+\epsilon\right)A_t
\right)
\right]
\end{equation}
と書かれる．ここで $\epsilon > 0$ は更新幅を制御するハイパーパラメータであり，$\mathrm{clip}(\cdot)$ は値を区間 $[1-\epsilon,1+\epsilon]$ に切り詰める演算である．新方策が旧方策から過度に離れると，利得の増加が頭打ちになるため，大きすぎる更新が抑制される．

PPO が実用上の標準手法とみなされるのは，信頼領域の思想を保ちつつ，実装が容易でミニバッチ最適化とも相性がよいからである．とくに連続行動空間，大規模関数近似，確率的方策の直接制御が必要な問題では，方策勾配法の重要性は大きい．価値ベース手法が価値推定を通じて行動を導くのに対し，方策勾配法は望ましい行動分布そのものを直接整形する．この違いが，現代の深層強化学習における両者の役割分担を定めている．

\subsection{発展的話題}

標準的な MDP，価値学習，方策勾配法は強力な枠組みであるが，現実世界の課題をそのまま捉えられるわけではない．実応用では，環境モデルは未知であり，観測は不完全であり，安全制約があり，データ収集が高価であることが多い．この節では，そのような制約のもとで基本理論がどのように拡張されるかを概観する．

\subsubsection{モデルベース強化学習}

モデルベース強化学習（model-based RL）が必要になるのは，サンプル効率を高めたいからである．モデルフリー手法は柔軟であるが，多数の試行を要する．実機ロボットや都市交通制御では，試行そのものが高価であり，環境モデルを全く使わない学習は非現実的であることが多い．

基本的な発想は，遷移モデル $\hat{P}$ と報酬モデル $\hat{r}$ をデータから学習し，その近似モデルの内部で計画やロールアウトを行うことである．モデルが十分正確なら，実環境での試行を減らしつつ方策改善ができる．一方で，モデル誤差は長期予測で蓄積しやすいため，短い計画区間に限定する，モデル不確実性を推定する，モデルフリー更新と併用するといった工夫が必要になる．

\subsubsection{オフライン強化学習}

オフライン強化学習（offline RL）が必要になるのは，新たな探索を現場で行えない状況があるからである．医療，金融，公共交通運用のような領域では，学習のために試しに危険な行動を取ることは許されない．そのため，過去に収集された固定データ集合
\begin{equation}
\mathcal{D} = \{(s,a,r,s')\}
\end{equation}
だけを用いて方策を学習する．

ここでの難しさは，データに含まれていない行動を価値関数が過大評価しやすい点にある．すなわち，分布外の行動に対する楽観主義が性能劣化を招く．このため，オフライン強化学習では，行動方策に近い範囲で保守的に学習する，価値推定を意図的に低めに補正する，行動分布そのものを制約する，といった工夫が中心課題となる．

\subsubsection{部分観測 MDP}

標準的 MDP では現在状態が将来予測に必要な情報を十分含むと仮定したが，現実には状態そのものが直接観測できないことが多い．部分観測 MDP（partially observable MDP; POMDP）は，この状況を扱うための拡張であり，エージェントは状態 $s_t$ の代わりに観測 $o_t$ だけを得る．

このとき単一の観測だけでは意思決定に必要な情報が不足することがある．たとえば需要の潜在状態や他主体の意図は直接見えない．そこで，過去観測から信念状態を更新して隠れ状態の分布を内部状態として扱う方法や，再帰型ネットワークによって観測履歴を潜在表現へ圧縮する方法が用いられる．POMDP は，マルコフ性を放棄するのではなく，観測から実質的なマルコフ表現を再構成する問題として理解できる．

\subsubsection{探索の高度化}

疎な報酬や長いホライズンを持つ問題では，単純な $\epsilon$-greedy では有益な経験に到達しにくい．このため，探索そのものを主題化する必要がある．代表的な考え方の一つが内発的動機づけ（intrinsic motivation）であり，新規性や予測誤差にもとづく探索ボーナスを外的報酬に加える．

たとえば修正報酬を
\begin{equation}
r_t' = r_t + \beta b_t
\end{equation}
と書けば，$\beta > 0$ は探索ボーナスの重み，$b_t$ は新規性や予測困難性を表す量である．ほかにも，不確実性の大きい行動を積極的に高く見る楽観主義や，事後分布から方策をサンプルする posterior sampling がある．探索は単なるランダム化ではなく，情報獲得をどう最適化するかという設計問題へ発展している．

\subsubsection{多エージェント強化学習}

多エージェント強化学習（multi-agent RL; MARL）が必要になるのは，環境が単一主体ではなく複数主体の相互作用で構成されるからである．交通流や配車市場では，各主体が学習することで他主体にとっての環境も変化する．その結果，各エージェントから見れば環境は非定常となり，単一エージェント MDP の仮定がそのままでは通用しにくい．

ここでの問題は，協調と競争をどう扱うかである．協調型では局所情報から全体効率を高める共同方策を構成する必要があり，競争型では相手の学習を見越した戦略的適応が求められる．実用的には，学習時には全体情報を利用し，実行時には各主体が局所観測だけで動く centralized training with decentralized execution が広く用いられる．

\subsubsection{模倣学習と逆強化学習}

現実の問題では，報酬関数を適切に設計すること自体が難しい．自動運転や運行制御では，「安全」「快適」「効率」といった複数の目的を単一の数値報酬に落とし込むのが容易ではない．このとき，人間や既存システムの行動データを利用する方法として，模倣学習（imitation learning）と逆強化学習（inverse RL）が重要になる．

模倣学習は専門家の行動系列から直接方策を学ぶ方法であり，探索コストを大きく削減できる．逆強化学習は，専門家行動の背後にある報酬関数を推定し，その推定報酬のもとで方策最適化を行う方法である．前者が「どう行動するか」を学ぶのに対し，後者は「何を目的としているか」を推定する．いずれも，標準的な MDP＋価値学習＋方策勾配だけでは扱いにくい現実的制約を補う拡張である．

\subsection{まとめ}

本章では，強化学習を環境との相互作用のもとで長期的な意思決定を最適化する理論として体系的に整理した．まず，教師あり学習や教師なし学習では十分に捉えられない遅延報酬，試行錯誤，探索と活用の問題から強化学習が必要になることを述べ，その歴史が動的計画法，TD 学習，Q 学習，深層関数近似，方策勾配法へと展開してきた流れを確認した．

続いて，マルコフ決定過程によって状態，行動，遷移確率，報酬，割引率，方策，価値関数を定義し，ベルマン方程式が強化学習全体の再帰的骨格であることを示した．多腕バンディット問題は，探索と活用の本質を状態遷移なしに抽出する縮約版として位置づけられ，動的計画法は環境モデル既知の場合の理論的基準点を与えた．そのうえで，モデル未知の状況では TD 学習と Q 学習が経験から価値を推定し，DQN が高次元状態へ拡張した．さらに，価値ベース手法の限界を受けて，方策勾配法，Actor-Critic，PPO が方策そのものを直接最適化する枠組みを与えることを見た．最後に，モデルベース強化学習，オフライン強化学習，POMDP，探索の高度化，多エージェント強化学習，模倣学習と逆強化学習を通じて，現実世界では未知性，部分観測，安全制約，データ収集制約が基本理論の拡張を要請することを確認した．以上を通じて，強化学習は「環境との相互作用のもとで価値関数または方策を学習し，長期的な意思決定を最適化する理論」として理解されるべきである．
